% Shared preamble for the EM-LLM Reader's Companion.
% Loaded once from main.tex; chapters never repeat \usepackage.

\usepackage[utf8]{inputenc}
\usepackage[T1]{fontenc}
\usepackage{lmodern}
\usepackage[a4paper,margin=1in]{geometry}
\usepackage{microtype}

\usepackage{amsmath}
\usepackage{amssymb}
\usepackage{amsthm}
\usepackage[ruled,linesnumbered]{algorithm2e}
\usepackage{array}
\usepackage{booktabs}
\usepackage{seqsplit}
\usepackage{xcolor}
\usepackage{enumitem}
\usepackage{csquotes}

\usepackage{listings}
\usepackage{tikz}
\usetikzlibrary{positioning,arrows.meta,calc,fit,backgrounds,decorations.pathreplacing,shapes.geometric}

\usepackage[backend=biber,style=numeric,sorting=none,maxbibnames=99]{biblatex}
\addbibresource{refs.bib}

\usepackage[colorlinks=true,linkcolor=black!70!blue,citecolor=black!50!blue,urlcolor=black!50!blue]{hyperref}
\usepackage[capitalize,noabbrev]{cleveref}

% Python listing style. The intent here is dense but readable code on the page,
% so we keep numbers on the left and avoid heavy frames that fight with the
% body text.
\definecolor{lstbg}{HTML}{F7F7F2}
\definecolor{lstkw}{HTML}{0B4F6C}
\definecolor{lststr}{HTML}{8B5E00}
\definecolor{lstcom}{HTML}{577590}

\lstdefinestyle{pythonstyle}{
    language=Python,
    backgroundcolor=\color{lstbg},
    basicstyle=\ttfamily\footnotesize,
    keywordstyle=\color{lstkw}\bfseries,
    stringstyle=\color{lststr},
    commentstyle=\color{lstcom}\itshape,
    numbers=left,
    numberstyle=\tiny\color{black!50},
    numbersep=8pt,
    showstringspaces=false,
    breaklines=true,
    breakatwhitespace=true,
    tabsize=4,
    frame=single,
    framerule=0.3pt,
    rulecolor=\color{black!30},
    captionpos=b,
    xleftmargin=10pt,
    xrightmargin=4pt,
    aboveskip=0.6em,
    belowskip=0.6em,
}
\lstset{style=pythonstyle}

% Shell snippets reuse the same frame but keep the default lexer.
\lstdefinestyle{shellstyle}{
    language=bash,
    backgroundcolor=\color{lstbg},
    basicstyle=\ttfamily\footnotesize,
    keywordstyle=\color{lstkw}\bfseries,
    commentstyle=\color{lstcom}\itshape,
    numbers=none,
    showstringspaces=false,
    breaklines=true,
    frame=single,
    framerule=0.3pt,
    rulecolor=\color{black!30},
    xleftmargin=10pt,
    xrightmargin=4pt,
}

% Convenience commands for frequently used objects.
% \seqsplit lets long identifiers and paths break across line ends without
% sliding into the right margin. Marked robust so the macros survive being
% used inside moving arguments (section titles, captions, TOC entries).
\DeclareRobustCommand{\file}[1]{\texttt{\seqsplit{#1}}}
\DeclareRobustCommand{\code}[1]{\texttt{\seqsplit{#1}}}
\newcommand{\param}[1]{\ensuremath{#1}}
\newcommand{\Surprise}{\mathcal{S}}
\newcommand{\Adj}{\mathbf{A}}
\newcommand{\Modularity}{Q}
\newcommand{\Conductance}{\phi}

\theoremstyle{definition}
\newtheorem{definition}{Definition}[section]
\newtheorem{example}{Example}[section]

\setlist{itemsep=2pt,topsep=4pt}
